\documentclass[11pt]{article}
\usepackage[final]{acl}
\usepackage{times}
\usepackage{latexsym}
\usepackage{booktabs}
\usepackage{multirow}
\usepackage[T1]{fontenc}
\usepackage[utf8]{inputenc}
\usepackage{microtype}
\usepackage{inconsolata}
\usepackage{xeCJK}
\setCJKmainfont{Songti SC}

\title{Chinese CEFR Labels Inherit an Ambiguity:\\
46\% of Texts Change Band with the Source Standard}

\author{Alvaro Serrano \\
  \textit{Independent researcher} \\
  \texttt{alvaro.serrano.gp@gmail.com} \\
  \texttt{orcid.org/0009-0006-4701-9026}}

\begin{document}
\maketitle

\begin{abstract}
Multilingual CEFR resources now span thirteen languages
\citep{imperial2025universalcefr}; Chinese is not among them. The obvious route
in is HSK, the Chinese proficiency scale, which is routinely mapped onto the
CEFR's six bands. That route is less well defined than it looks. ``HSK 3.0''
names two official documents --- the GF0025-2021 national grading standard and
the examination syllabus in force since July 2026 --- which assign different
levels to 41.5\% of the vocabulary they share \citep{serrano2026whichhsk}. We
show the ambiguity propagates: grading 102 graded readers against one document
rather than the other moves \textbf{46.1\%} of them into a different CEFR band
under the usual one-to-one correspondence, and \textbf{44.1\%} under a
compressed correspondence that assumes HSK~6~$\approx$~B2. It replicates on a
disjoint 30-text held-out split (46.7\%, 43.3\%). The movement is not noise:
\textbf{42 of 47} texts move \emph{upward}, and on the held-out split all 14
do. Choosing the wrong document does not add error, it shifts the distribution.
We argue that any Chinese entry in a CEFR corpus must record which Chinese
standard its labels derive from, and that no correspondence table can repair
the problem, because the ambiguity sits upstream of the mapping.
\end{abstract}

\section{The gap, and why it is awkward to fill}

\citet{imperial2025universalcefr} assemble 494{,}491 CEFR-labelled texts across
thirteen languages and standardise them into one schema. Chinese is absent.
This is not an oversight of scale --- Chinese is among the most-studied second
languages in the world --- but a consequence of the inclusion criteria, which
require a permissive licence, CEFR labels produced or validated by domain
experts, and human authorship.

The natural way to produce Chinese CEFR labels at scale is indirect: grade text
against HSK, then map HSK levels onto CEFR bands. Learner-facing material does
this constantly. The mapping is not official --- China's own standard declines
to align itself to the CEFR --- but it is ubiquitous, and it is the implicit
basis on which Chinese material is described as ``B1'' or ``C1'' in practice.

This note reports an obstacle that sits before the mapping, and that a better
mapping cannot fix.

\section{Two documents named HSK 3.0}
\label{sec:two}

The name ``HSK 3.0'' is applied to two distinct official publications:

\begin{itemize}
\item \textbf{《国际中文教育中文水平等级标准》} (GF0025-2021), the national
grading standard, in force since 1 July 2021. Nine levels, 11{,}092 words,
3{,}000 characters.
\item \textbf{新版HSK考试大纲}, the examination syllabus, published November
2025 and in force since July 2026. 10{,}896 words, 3{,}088 recognition
characters, and a separate 1{,}200-character writing inventory.
\end{itemize}

\citet{serrano2026whichhsk} compares them directly: of the 9{,}674 words the
two share, 41.5\% sit at different levels; of the 2{,}945 shared characters,
40.7\% do. Grading a corpus of authentic graded readers against one rather than
the other changes the assigned HSK level of 48.0\% of texts. Tools and
reference sites that report ``HSK 3.0'' levels overwhelmingly do not say which
document they used.

The question this note asks is narrow: does that disagreement survive
projection onto the CEFR's coarser six-band scale, or does coarsening wash it
out?

\section{Method}

\paragraph{Corpus.} 102 graded readers spanning six difficulty shelves, with
per-word segmentation authored rather than produced by a segmenter, plus a
disjoint 30-text held-out split. Both are released CC~BY~4.0 with a datasheet.
We grade the authored segmentation directly; re-tokenising the raw strings with
a segmenter yields materially different numbers and is not comparable to the
published figures.

\paragraph{Grading.} Each text is graded twice with the same open
implementation \citep{serrano2026whichhsk}, once against GF0025-2021 and once
against the 2025 syllabus. A text's level is the lowest level whose cumulative
character inventory covers 95\% of its running characters, the minimal coverage
threshold from \citet{laufer2010lexical}. Nothing about the procedure differs
between the two runs except the inventory.

\paragraph{Correspondences.} We use two, and report every figure under both.
The \textbf{naive} correspondence is the one-to-one reading ubiquitous in
learner material: HSK~$n$ maps to the $n$th CEFR band, with HSK 3.0's combined
levels 7--9 mapped to C2. The \textbf{compressed} correspondence encodes the
objection practitioners raise --- that vocabulary coverage overstates
communicative range, so HSK~6 lands nearer B2 than C2 --- by mapping
$\{1,2\}\!\to\!$A1, $3\!\to\!$A2, $4\!\to\!$B1, $\{5,6\}\!\to\!$B2,
$\{7,8\}\!\to\!$C1, $9\!\to\!$C2.

Neither is authoritative. That is the point of using both: a finding that holds
under both does not rest on the choice.

\section{Results}

\begin{table}[t]
\centering
\small
\begin{tabular}{lrr}
\toprule
& \textbf{Readers} & \textbf{Held-out} \\
& ($n$=102) & ($n$=30) \\
\midrule
HSK level changes        & 49 (48.0\%) & 14 (46.7\%) \\
\addlinespace
CEFR band, naive         & 47 (46.1\%) & 14 (46.7\%) \\
\quad upward             & 42          & 14 \\
\quad downward           & 5           & 0 \\
\addlinespace
CEFR band, compressed    & 45 (44.1\%) & 13 (43.3\%) \\
\quad upward             & 42          & 13 \\
\quad downward           & 3           & 0 \\
\bottomrule
\end{tabular}
\caption{Texts whose band differs depending on which HSK document was used.
``Upward'' means the 2025 syllabus placed the text higher than the 2021
standard did.}
\label{tab:main}
\end{table}

\begin{table}[t]
\centering
\small
\begin{tabular}{lrrrrrr}
\toprule
& \textbf{A1} & \textbf{A2} & \textbf{B1} & \textbf{B2} & \textbf{C1} & \textbf{C2} \\
\midrule
2021 basis & 3 & 28 & 32 & 18 & 14 & 7 \\
2025 basis & 5 & 10 & 29 & 32 & 19 & 7 \\
\bottomrule
\end{tabular}
\caption{Band distribution of the same 102 texts under the naive
correspondence. The A2 population falls by 64\% and the B2 population rises by
78\%; the corpus is not relabelled at random, it is relabelled harder.}
\label{tab:dist}
\end{table}

Table~\ref{tab:main} gives the result. Coarsening to six bands does not wash
the disagreement out: 46.1\% of texts change band under the naive
correspondence, against 48.0\% that change HSK level. Compressing the
correspondence --- which merges adjacent HSK levels and so should absorb
single-level shifts --- reduces it only to 44.1\%. The pattern replicates on
the held-out split at 46.7\% and 43.3\%.

The direction is the more consequential finding. Of the 47 texts that move
under the naive correspondence, 42 move upward; on the held-out split, all 14
do. Table~\ref{tab:dist} shows the effect on the corpus as a whole: the A2
population falls from 28 to 10 while B2 rises from 18 to 32. A researcher who
picks the wrong document does not acquire noisy labels. They acquire a corpus
that is systematically mis-levelled in one direction, which is the failure mode
least likely to be caught by inspection and most likely to bias a model trained
on it.

\section{What this implies for multilingual CEFR resources}

\paragraph{The provenance field is not optional.} The
\citet{imperial2025universalcefr} schema records source dataset, language,
granularity, production category and licence. For a Chinese entry this is
insufficient: two entries could share every one of those values, be produced by
the same procedure over the same texts, and still disagree on the band of
nearly half of them. Chinese needs the source standard, with its version,
recorded as data. We would expect the same to apply wherever a language's CEFR
labels are derived through a national framework rather than assigned directly.

\paragraph{No correspondence table fixes it.} It is tempting to read this as an
argument for a better HSK--CEFR mapping. It is not. Both correspondences we
test are applied identically to both gradings; the disagreement is generated
before the mapping is consulted and passes through it. A third correspondence
would behave the same way.

\paragraph{Derived labels are not expert labels.} We are explicit that the
labels here are algorithmic projections of a coverage threshold, not
expert-assigned CEFR levels, and so do not meet the gold-standard criterion of
\citet{imperial2025universalcefr}. The contribution is not a proposed dataset
entry. It is a constraint that a future entry will have to satisfy.

\section{Limitations}

The corpus is LLM-assisted, disclosed in its datasheet, and pedagogically
constrained; it is not a sample of naturally occurring Chinese and would not
meet the human-authorship criterion of \citet{imperial2025universalcefr}. Its
levels are ours, not experts'. Coverage is necessary but not sufficient for
comprehension: grammar, register and world knowledge are unmodelled, and the
95\% threshold is imported from word-level English research rather than
established for L2 Chinese characters. Neither correspondence we use is
official, and we make no claim that either is correct --- only that the finding
does not depend on which is chosen. Finally, the 2025 syllabus takes effect in
July 2026, so its practical consequences are not yet observable in learner
outcomes.

\section{Availability}

The grading implementation is MIT licensed and installable as
\texttt{pip install hsk30}; the corpus and held-out split are CC~BY~4.0. Every
figure in this note is regenerated by \texttt{scripts/cefr\_mapping.py} in the
same repository, which prints the table above from the corpus and the two
inventories.

\section*{Ethics and data statement}

The corpus contains no personal data and describes invented characters in
everyday situations; its LLM-assisted authorship is disclosed in its datasheet
and again in the limitations above. Word and character lists derive from
published national standards. No human subjects were involved. An intended
application is a Chinese-learning website operated by the author.

\bibliography{refs}

\end{document}
